% !TEX root = ../print.tex
% Draft \S11 --- The scale evolution in generator form
%
% Source: draft/spatial-hemigroup-scale-space.md, \S11 (lines 665-699).
%
% CHANGED (prop:scale-evolution): the draft leaves the signal class open, with a check
% asking for the class of f on which the pointwise identity holds. The node fixes it as the
% Schwartz class, which is the smallest choice that makes every step of the proof
% unconditional and is enough for every later use (thm:non-enhancement, prop:jet-regularity).
% Whether the evolution extends to a larger domain -- W^{2,1} for the Gaussian part, the full
% \Levy{} generator domain in general -- is a question for the paper on the infinitesimal theory
% and is not answered here.
%
% CHANGED (prop:corner-generators clause 4): the draft carries a check on the constants and
% on whether the closed-form Student-t symbol is worth stating. The node states the Thorin
% mixture with the constants worked out from prop:scale-evolution, and does NOT state the
% closed form, which is not used anywhere.
%
% The full infinitesimal theory -- the scale-Cauchy problem, the identification of the
% Choquet coordinates with the atomic generators, and Phillips' formula as the infinitesimal
% form of the bridge -- is deferred to the paper on the infinitesimal theory; the drafted
% statements are parked in deferred/scale-evolution.md and are not part of any release.

\chapter{The scale evolution in generator form}
\label{ch:generator}

The classification is stated on the kernels; the chapters that follow use it on the evolution
across scale, in the form the subject was first posed in. This chapter states that form and the
generators of the corners, which is all the present article needs.

% additive: the signal class, fixed here rather than left open; see the header note.
% CHANGED (review 2026-09-09, Q9): the annotation now names the three Mathlib declarations that
% carry the definition's three ancillary properties. The \lean tag names the class alone, which
% is right for a definition node, and the annotation is where the provenance of the properties
% belongs -- naming them there documents them without creating a proof obligation.
% CHANGED (2026-09-19, article mirror of module B, R174): the annotation said the Schwartz class is
% "the smallest choice" that makes prop:scale-evolution unconditional, an unproved superlative (the
% register pass's item 16; the header note above uses the same word and is a comment, not the text
% of record). It now says what is true: a choice, the simplest this article found, with no
% minimality claimed. Annotation only; the definition, the tag and the status are untouched.
\begin{definition}[the signal class]\label{def:signal-class}
\lean{SpatialLine.SignalClass}\leanok
$\mathcal{D} := \Sch$ is the Schwartz class on the line: smooth functions all of whose
derivatives decay faster than every polynomial. It is contained in $\Lone$, is stable under
differentiation, and its Fourier image is itself.
\statusT\quad\emph{(A choice that makes every step of \ref{prop:scale-evolution}
unconditional, and the simplest this article found; that no smaller class does is not claimed,
and no other class is compared with it here. Enlarging it is a question for the paper on the
infinitesimal theory. The three
ancillary properties are standard and are carried in the formalisation by Mathlib rather than
by declarations of ours: \texttt{SchwartzMap.integrable} for the inclusion in $\Lone$,
\texttt{SchwartzMap.derivCLM} for stability under differentiation, and
\texttt{SchwartzMap.fourierTransformCLE} for the Fourier image, with stability under convolution
with a finite measure following from the first.)}
\end{definition}

% draft: Proposition 11.1
% CHANGED (proof of record 2026-09-10): the proof of eq:evolution-signal is rewritten to the
% machine-checked route. The printed proof read the identity as "the inverse transform of
% eq:evolution-fourier", which presupposes that the generator may be applied under the inverse
% transform at u(t,.). It may not: A_t is a Fourier multiplier on the Schwartz class and u(t,.)
% is not Schwartz, since it inherits the kernel's tails. The commutation A_t(nu * f) =
% nu * (A_t f), which is the step that repairs this, is now stated in the proof, and the
% inversion is stated once as an identity for a symmetric probability measure and used twice.
% The dominating function is named, and so is the decay of the pairing that makes it integrable
% -- one power more than the symbol's growth consumes.
% twin: hcs prop:scale-evolution and def:phillips-generator
% (Hemigroup.SelfDecomposableExponent.phillipsGenerator) --- the causal generator has the
% memory kernel where this one has symmetric second differences.
\begin{proposition}[the scale evolution in generator form]\label{prop:scale-evolution}
\uses{thm:main-characterization,lem:selfdecomposable-exponents,def:signal-class,lem:quadratic-growth}
\lean{SpatialLine.scale_evolution_fourier, SpatialLine.scale_evolution_absconv,
SpatialLine.scale_evolution_multiplier, SpatialLine.scale_evolution_signal,
SpatialLine.scale_evolution_conjugation}\leanok
\notes{covariant-memory-evolution}
Let $F$ be admissible with data $(a,k)$ and $\varpi = -dk$, and let $u(t,\cdot) :=
\mu_{0,t} * f$ be the scale space of $f \in \mathcal{D}$ in the canonical gauge. Then on the
Fourier side
\begin{equation}\label{eq:evolution-fourier}
  \partial_t\hat u(t,\omega) = -\frac1t\,B(t\omega)\,\hat u(t,\omega),
  \qquad
  B(\omega) = \omega F'(\omega)
   = 2a\,\omega^2 + \int_{(0,\infty)}\bigl(1 - \cos\omega v\bigr)\,\varpi(dv),
\end{equation}
for every $\omega$ and every $t > 0$, and on the signal side
\begin{equation}\label{eq:evolution-signal}
  \partial_t u(t,x) = \mathcal{A}_t u(t,x),
  \qquad
  \mathcal{A}_t g(x) := 2at\,g''(x)
   - \frac1t\int_{(0,\infty)}\Bigl[g(x) - \tfrac12\bigl(g(x-tv) + g(x+tv)\bigr)\Bigr]\,\varpi(dv),
\end{equation}
the integral converging absolutely for $g \in \mathcal{D}$. The generators at different scales
are conjugate, $\mathcal{A}_t = t^{-1}\Dil{t}\,\mathcal{A}_1\,\Dil{t}^{-1}$, with
$\mathcal{A}_1 = -B(D)$ and $D = -i\partial_x$; in log-scale $\theta = \log t$ the
evolution reads $\partial_\theta u = \Dil{t}\mathcal{A}_1\Dil{t}^{-1}u$.
\statusT\quad\emph{(The operator $\mathcal{A}_1 = -B(D)$ is the generator of the symmetric
\Levy{} process with Gaussian part $2a$ and folded jump measure $\varpi$. That reading stood in the
conjugation clause until 2026-09-19 (R174) and is here instead: it names an object this article
never constructs, no declaration of the tag reads it, and as printed it attached the words ``the
generator of'' to $D = -i\partial_x$, which is not one. The two forms are the two readings of
\ref{rem:displacement-dictionary} as operators: the Gaussian coefficient is a Laplacian, and
the catalogue is a superposition of symmetric second differences over displacement sizes, at
rate $\varpi$ and dilated to the current scale. The factor $t^{-1}$ says that a unit of
log-scale, not of scale, is the natural clock. The signal-side identity is a pointwise
derivative in the scale at a fixed position, not a derivative in $\Lone$ or in any ambient
space; the larger domain that \ref{def:signal-class} defers is a question for the paper on
the infinitesimal theory. \texttt{\#print axioms} on each of the five declarations is Lean
core. \emph{Both derivative clauses name their derivative}: they are stated with
\texttt{HasDerivAt}, which asserts existence and value together, so neither reads the default
value of a derivative that may not exist. The generator on the right of the signal-side clause
contains a second space derivative of the scale space, and that is not a default either: the
scale space of a signal in $\mathcal{D}$ is \emph{proved} twice differentiable, by
differentiation under the convolution integral
(\texttt{hasDerivAt\_mconv}, \texttt{deriv\_mconv\_signal},
\texttt{iteratedDeriv\_two\_mconv}), and the checked proof of the signal-side clause routes
through that chain rather than around it. This closes the review's row R6, which had priced
the missing guard as a defect; the guard exists and is the convolution's own smoothing.)}
\end{proposition}

\begin{proof}
\leanok
$\hat u(t,\omega) = e^{-F(t\omega)}\hat f(\omega)$, and $F$ is $C^1$ off the origin with
$\omega F'(\omega) = B(\omega)$ by \ref{lem:selfdecomposable-exponents} and \ref{eq:symbol};
differentiating in $t$ gives \ref{eq:evolution-fourier}, since
$\partial_t F(t\omega) = \omega F'(t\omega) = t^{-1}B(t\omega)$.

For \ref{eq:evolution-signal}: $B(t\omega) = 2at^2\omega^2 + \int(1 - \cos t\omega
v)\,\varpi(dv)$, where $\omega^2$ is the symbol of $-\partial_x^2$ and $1 - \cos(\omega\,tv)$
is the symbol of $g \mapsto g - \tfrac12(\Tr{tv} + \Tr{-tv})g$. So the Fourier multiplier
$B(tD)$ acts on $g \in \mathcal{D}$ as
$-2at^2g'' + \int[g - \tfrac12(g(\cdot - tv) + g(\cdot + tv))]\,\varpi(dv)$, the integrand
being bounded by $\min(2\|g\|_\infty, \tfrac12t^2v^2\|g''\|_\infty)$ and hence
$\varpi$-integrable, since $\int (1 \wedge v^2)\,\varpi < \infty$. Thus
$\mathcal{A}_t = -t^{-1}B(tD)$ on $\mathcal{D}$, which is \ref{eq:evolution-signal} read as a
Fourier multiplier.

That $\partial_t u$ equals $\mathcal{A}_tu$ pointwise is the inverse transform of
\ref{eq:evolution-fourier}, and the inversion is used twice. Write
$\langle\kappa_\varphi, h\rangle = \int_\RR h(z)\cos(\omega z - \varphi)\,dz$ for the pairing of
$h$ with the character of frequency $\omega$ and phase $\varphi$. For a symmetric probability
measure $\nu$ and an $h$ that is continuous and integrable with integrable transform,
\[
  (\nu * h)(x) \;=\; \frac{1}{2\pi}\int_\RR \hat\nu(\omega)\,
    \langle\kappa_{\omega x}, h\rangle\,d\omega ,
\]
by Fourier inversion for $h$ together with one exchange of integrals against the finite $\nu$;
the symmetry of $\nu$ is what turns the $\nu$-average of $\cos(\omega(x-y))$ into
$\hat\nu(\omega)\cos(\omega x)$, and no other step of the proof uses it.

Read at $h = f$ and $\nu = \mu_{0,s}$, this represents the scale space as
$u(s,x) = (2\pi)^{-1}\int e^{-F(s\omega)}\langle\kappa_{\omega x},f\rangle\,d\omega$.
Differentiating under that integral at $s = t$ is licensed on $s \in (t/2, 3t/2)$ by the
dominating function $s^{-1}B(s\omega)\,|\langle\kappa_{\omega x},f\rangle|$, which is integrable
because $B(s\omega) \le C(1 + s^2\omega^2)$ by \ref{lem:quadratic-growth} while
$|\langle\kappa_\varphi, f\rangle| \le (\|f\|_1 + \|f''\|_1)/(1 + \omega^2)$ uniformly in the
phase --- the two crude bounds $|\langle\kappa_\varphi,f\rangle| \le \|f\|_1$ and
$\omega^2|\langle\kappa_\varphi,f\rangle| = |\langle\kappa_\varphi,f''\rangle| \le \|f''\|_1$,
combined --- and iterating that estimate on $f''$ improves the decay to
$(1 + \omega^2)^{-2}$, which is one power more than the symbol's growth consumes. So
$\partial_tu(t,x)
 = -(2\pi)^{-1}\int t^{-1}B(t\omega)e^{-F(t\omega)}\langle\kappa_{\omega x},f\rangle\,d\omega$.

Read at $h = \mathcal{A}_tf$ and $\nu = \mu_{0,t}$, it represents the right-hand side --- but
only after the generator has been moved off $u(t,\cdot)$, which is \emph{not} in $\mathcal{D}$:
the scale space inherits the kernel's tails, and the Student-t and stable kernels have
polynomial ones, so the multiplier identity above does not
apply to it. Translation invariance supplies the step: $\mathcal{A}_t$ commutes with convolution
by a finite measure, $\mathcal{A}_t(\nu * f) = \nu * (\mathcal{A}_tf)$, the exchange of the
$\varpi$-integral with the $\nu$-integral being licensed by the same bound
$\min(2\|f\|_\infty, \tfrac12t^2v^2\|f''\|_\infty)$ that makes the generator's own integral
converge. Hence $\mathcal{A}_tu(t,\cdot) = \mu_{0,t} * (\mathcal{A}_tf)$, whose argument is
continuous and integrable, and
$\langle\kappa_{\omega x},\mathcal{A}_tf\rangle = -t^{-1}B(t\omega)\langle\kappa_{\omega x},f\rangle$
by the multiplier identity at the scale $t$. The two frequency integrands agree, which is
\ref{eq:evolution-signal}.

Conjugation: $\Dil{t}$ has symbol map $\omega \mapsto t\omega$, so
$\Dil{t}B(D)\Dil{t}^{-1} = B(tD)$; and $\partial_\theta = t\partial_t$.
\end{proof}

% draft: Proposition 11.2
% twin: hcs has no single node for the corner generators; the causal article gives its
% generator once, in def:phillips-generator, and reads the corners off it in prose.
% SKELETON NOTE (phase B, 2026-09-09): clause (2) is NOT typed as the operator identity
% A_t = -alpha t^{alpha-1}(-d_x^2)^{alpha/2}: Mathlib has no fractional Laplacian. What
% Skeleton.corner_generator_stable states is the pair that clause is about and that the node's
% own proof computes -- the symbol B(omega) = alpha |omega|^alpha and the jump measure
% varpi(dv) = C_alpha^{-1} alpha v^{-alpha-1} dv -- which together determine the operator.
% SKELETON.md question Q12.
% CHANGED (review 2026-09-09, Q12): the annotation now says that clause (2) is formalised at the
% symbol and the jump measure, and why that is a faithful reading of the clause rather than a
% gap -- the multiplier identity of prop:scale-evolution turns the pair into the operator, and
% both are typed. No Lean change and no statement change.
% CHANGED (proof 2026-09-10): the class of g in the FORMALISATION of clause (4) is narrowed to
% the class of prop:scale-evolution's convergence clause -- g twice continuously differentiable
% with g and g'' bounded -- in place of the review's "measurable and bounded" (R18). Reason: the
% two sides are Bochner integrals and for a general Thorin member varpi is infinite near the
% origin, so for a merely bounded measurable g the varpi-integral of the second difference need
% not converge and the two sides junk independently; the Matern clause is safe as it stands
% because there varpi is finite. The narrowing is to the class the clause is about, and it is
% what the proof consumes: that convergence clause IS the integrability the Fubini needs. The
% blueprint statement, which is an operator identity and quantifies over no g, is unchanged.
% CHANGED (proof of record 2026-09-10): the derivation of the Choquet measure in the proof of
% (3) and (4) is rewritten to the checked route -- the measure is identified from its tails, not
% obtained by differentiating the profile -- and the exchange of the two integrals is named.
% CHANGED (2026-09-11, the author's decision on the post-sync recommendation, R116): draft
% Proposition 11.2(2) glossed its equation, inside the statement, as "the alpha-scale-space
% equation of [Duits et al. 2004] (their alpha being half of this one)", which this node did not
% carry. The connection is kept and moved out of the statement: it is now a sentence of the status
% note, at the pinned citekey duits2004axioms with the anchor read by the librarian from the held
% copy (p. 288, section 6.1 and eq. (66), generator -(-Delta)^alpha with 0 < alpha <= 1; p. 293,
% Fourier multiplier exp(-s |omega|^{2 alpha})). No clause and no Lean changed.
% CHANGED (2026-09-19, article mirror of module B, R174): two register repairs, from the module B
% register pass's items 7 and 13.
% (a) CLAUSE (4) NAMES ITS CLASS. The clause was an operator identity quantifying over no g, while
% its annotation and the paper both read it on the C^2 class with g and g'' bounded -- and both
% described that class as "the class of prop:scale-evolution's convergence clause", which is stated
% for Schwartz g. The clause now carries the class, which is a NARROWING (an identity at every g of
% a named class, where the print left the domain open), and the annotation says what the class is
% and that it is not the one it was attributed to. The Lean is unaffected: the narrowing of the
% FORMALISATION to this class was made on 2026-09-10 (the wave 6 marker above), so the statement now
% matches its declarations.
% (b) Clause (3)'s em-dash aside becomes a second sentence, with the same content: the operator is
% bounded, and it is the generator of a compound Poisson process with the same jump law and rate.
% CHANGED 2026-09-19 (R182, second external review round, referee C findings 8 and 9): TWO
% REPAIRS, both in the safe direction. (a) Clause (4)'s Student-t sentence was unconditional and
% mis-referenced: that the Student-t family is a Thorin member at all is conditional on ledger A18
% (rem:student-ggc), as prop:student-t(3) and lem:student-subordinated have said since R160, and the
% Bessel Thorin measure is exhibited in that remark and not in prop:thorin-subclass, whose clause
% (3) is the bridge. The sentence now carries the hypothesis and points at eq:student-ggc for the
% measure, with prop:thorin-subclass(3) for the image. (b) The annotation's reason for the class of
% test functions was false as stated -- "for a Thorin member other than the Matern one the Choquet
% measure is infinite near the origin" -- since varpi((0,infty)) = U((0,infty)) = k(0+) is finite
% for every Thorin member with finite U, which the proof below says in as many words ("as it does
% whenever U is finite"). It is now stated at k(0+) = infinity. No clause of the statement and no
% declaration is affected.
\begin{proposition}[the generators of the corners]\label{prop:corner-generators}
\uses{prop:scale-evolution,prop:matern-exponent,prop:student-t,prop:stable-family,prop:thorin-subclass,lem:cin-rays}
\lean{SpatialLine.corner_generator_gaussian, SpatialLine.corner_generator_stable,
SpatialLine.corner_generator_matern, SpatialLine.corner_generator_thorin,
SpatialLine.corner_generator_cin}\leanok
In the canonical gauge, and with $\Lap_\theta$ denoting convolution by the Laplace kernel
$\tfrac{1}{2\theta}e^{-|x|/\theta}$ of range $\theta$:
\begin{enumerate}
\item \emph{Gaussian:} $\mathcal{A}_t = 2at\,\partial_x^2$, the heat equation
$\partial_v u = \partial_x^2u$ in the variance gauge $v = at^2$.
\item \emph{Symmetric stable,} $F = |\omega|^\alpha$ with $0 < \alpha < 2$:
$\mathcal{A}_t = -\alpha\,t^{\alpha-1}(-\partial_x^2)^{\alpha/2}$, the fractional Laplacian;
in $v = t^\alpha$ the equation is $\partial_v u = -(-\partial_x^2)^{\alpha/2}u$, and the jump
measure is $\varpi(dv) = C_\alpha^{-1}\alpha\,v^{-\alpha-1}dv$.
\item \emph{\Matern{},} $F = \gamma\log(1+\omega^2)$:
$\mathcal{A}_t = \tfrac{2\gamma}{t}\,(\Lap_t - I)$, a \emph{bounded} operator of norm at most
$4\gamma/t$. It is the generator of a compound Poisson process, whose jumps are
Laplace-distributed of range $t$ and arrive at the rate $2\gamma/t$.
\item \emph{Thorin members,} with Thorin measure $U$ as in \ref{eq:thorin} and $a = 0$:
\[
  \mathcal{A}_t = \frac1t\int_{(0,\infty)}\bigl(\Lap_{t/\theta} - I\bigr)\,U(d\theta),
\]
the Thorin mixture of \Matern{} generators, as an identity at every $g$ that is twice continuously
differentiable with $g$ and $g''$ bounded. Under the hypothesis of \ref{rem:student-ggc} the
Student-t family is the case of the Bessel Thorin measure: $U$ is twice the image of the causal
$U_{\mathrm I}$ of \ref{eq:student-ggc} under $\theta \mapsto \sqrt{2\theta}$, by
\ref{prop:thorin-subclass}(3).
\item \emph{$\Cin$ ray of step $\tau$:} $\mathcal{A}_t = \tfrac1t\,\delta^2_{t\tau}$, a single
symmetric second difference, where
$\delta^2_h g := \tfrac12\bigl(g(\cdot - h) + g(\cdot + h)\bigr) - g$.
\end{enumerate}
\statusT\quad\emph{(Four kinds of generator: a differential operator, a fractional differential
operator, a bounded integral operator, and a mixture of bounded integral operators, with the
extreme rays carrying finite-difference generators. Clause (2)'s equation in $v = t^\alpha$ is the
$\alpha$-scale-space equation of \sscite{duits2004axioms} (p.~288, (66)), whose generator
$-(-\Delta)^{\alpha'}$, $0 < \alpha' \le 1$, is this one at $\alpha' = \alpha/2$: their exponent
is half of the stable index used here. Of the four, the first is the one whose generator is a
local operator, and which member of the cone that singles out is left for further work
(2026-09-19, R195: the reference into the locality chapter is replaced by this sentence, module C
being unpublished);
% TODO(module C): cite thm:scale-locality here once module C is released.
the boundedness of the third is the infinitesimal
face of \ref{thm:matern}. Clause (2) is formalised at the symbol and the jump measure rather
than at the operator: Mathlib has no fractional Laplacian, and the pair
$\bigl(B(\omega) = \alpha|\omega|^\alpha,\ \varpi\bigr)$ is what the clause's own proof computes
and what \emph{defines} the operator here, through the multiplier identity of
\ref{prop:scale-evolution}. The operator form is the reading of that pair, not a further
claim. Clauses (3) and (4) are formalised as pointwise identities in $x$, at a fixed scale, with
$\Lap_\theta$ written as convolution by the Laplace law of range $\theta$. The class named in
clause (4) --- $g$ twice continuously differentiable with $g$ and $g''$ bounded --- is the class on
which the bound in the proof of \ref{prop:scale-evolution} makes the generator's integral converge;
it is larger than $\mathcal{D}$, which is the class that proposition's convergence clause is stated
on, and the two were conflated here and in the proof below until 2026-09-19 (R174). The class is
needed because for a Thorin member with $U\bigl((0,\infty)\bigr) = k(0{+}) = \infty$ --- the
symmetric stable members among them ---
the Choquet measure is infinite near the origin and boundedness alone does not make
the generator's integral converge. For a member with $U$ finite, the \Matern{} one among them,
boundedness of $g$ suffices, as the proof records; the sentence here said ``other than the
\Matern{} one'', which is false for every Thorin member with finite $U$ --- two atoms, or any
finite quadrature of \ref{prop:thorin-machine} --- and is corrected on 2026-09-19 (R182).)}
\end{proposition}

\begin{proof}
\leanok
(1) is \ref{prop:scale-evolution} with $\varpi = 0$.

(2) $B(\omega) = \omega F'(\omega) = \alpha|\omega|^\alpha$, so
$-t^{-1}B(t\omega) = -\alpha t^{\alpha-1}|\omega|^\alpha$, and $|\omega|^\alpha$ is the symbol
of $(-\partial_x^2)^{\alpha/2}$; the change of variable $v = t^\alpha$ gives
$\partial_v = (\alpha t^{\alpha-1})^{-1}\partial_t$. The jump measure is $-dk$ for
$k(x) = C_\alpha^{-1}x^{-\alpha}$ (\ref{prop:stable-family}).

(3) and (4) are one computation. For $\theta > 0$,
$\int_0^\infty \bigl[g - \tfrac12(g(\cdot - v) + g(\cdot + v))\bigr]\,\theta e^{-\theta v}\,dv
= g - \Lap_{1/\theta}g$, because $\int_0^\infty \theta e^{-\theta v}dv = 1$ and
$\int_0^\infty \tfrac12(g(x-v)+g(x+v))\,\theta e^{-\theta v}\,dv
= \int_\RR g(x-y)\,\tfrac12\theta e^{-\theta|y|}\,dy = (\Lap_{1/\theta}g)(x)$.

A Thorin member with $a = 0$ has $k(x) = \int e^{-\theta x}U(d\theta)$
(\ref{prop:thorin-subclass}), and its Choquet measure is the $U$-mixture of exponential laws,
$\varpi = \int \theta e^{-\theta v}\,U(d\theta)\,dv$ --- equivalently, the image of the product
of $U$ with the standard exponential law under $(\theta,u) \mapsto u/\theta$. This is an
identification of
measures and not a differentiation of the profile: $\varpi$ is only given by its tails, and
almost everywhere at that, so what is checked is that the mixture has the right ones,
$\int_y^\infty\!\int \theta e^{-\theta v}U(d\theta)\,dv = \int e^{-\theta y}U(d\theta) = k(y)$
by Tonelli, after which two folded measures with the same tails on $(0,\infty)$ agree. Writing
the mixture as an image measure rather than as a density is what makes the exchange below the
ordinary Fubini on a product.

Substituting into \ref{eq:evolution-signal} and exchanging the $U$-integral with the
$\varpi$-integral --- licensed by the absolute convergence of \ref{prop:scale-evolution} for $g$
bounded with bounded second derivative --- and dilating $v \mapsto tv$, which sends
$\Lap_{1/\theta}$ to $\Lap_{t/\theta}$, gives
$\mathcal{A}_tg = \tfrac1t\int (\Lap_{t/\theta} - I)g\,U(d\theta)$, which is (4). For (3) the
\Matern{} member has $U = 2\gamma\delta_1$ by \ref{prop:matern-exponent}(1) and
\ref{prop:thorin-subclass}, so $\mathcal{A}_t = \tfrac{2\gamma}{t}(\Lap_t - I)$; the norm bound
is $\|\Lap_t\| \le 1$ on $\Lone$ and on $L^\infty$. There $\varpi = 2\gamma e^{-v}dv$ is finite
and boundedness of $g$ alone suffices, as it does whenever $U$ is finite, since
$\varpi((0,\infty)) = U((0,\infty)) = k(0{+})$; a Thorin member with $k(0{+}) = \infty$ has
$\varpi$ infinite near the origin, which is why clause (4) names the class it does.

(5) is \ref{prop:scale-evolution} with $a = 0$ and $\varpi = \delta_\tau$, the profile being
$\mathbf 1_{(0,\tau)}$ (\ref{lem:cin-rays}(1)).
\end{proof}
